
\documentclass{article}
\usepackage{colortbl}
\usepackage{makecell}
\usepackage{multirow}
\usepackage{supertabular}

\begin{document}

\newcounter{utterance}

\centering \large Interaction Transcript for game `function\_detective', experiment `list\_passive\_examples\_oneshot', episode 118 with mistral{-}large{-}3{-}azure.
\vspace{24pt}

{ \footnotesize  \setcounter{utterance}{1}
\setlength{\tabcolsep}{0pt}
\begin{supertabular}{c@{$\;$}|p{.15\linewidth}@{}p{.15\linewidth}p{.15\linewidth}p{.15\linewidth}p{.15\linewidth}p{.15\linewidth}}
    \# & \multicolumn{2}{c}{Player} && \multicolumn{2}{c}{Game Master} \\
    \hline

    \theutterance \stepcounter{utterance}  
    & & & \multicolumn{4}{p{0.6\linewidth}}{
        \cellcolor[rgb]{0.9,0.9,0.9}{
            \makecell[{{p{\linewidth}}}]{
                \texttt{\tiny{[P1$\langle$GM]}}
                \texttt{I have a hidden Python function for you to figure out.} \\
\\ 
\texttt{Your job is to infer what it does and then submit a correct Python implementation.} \\
\\ 
\texttt{Function signature:} \\
\texttt{(items: List{[}int{]}) {-}> int} \\
\\ 
\texttt{You have 1 turns total.} \\
\\ 
\texttt{I have already given you a fixed list of valid input{-}output examples. Solve using those examples only.} \\
\\ 
\texttt{You must make your best SOLVE attempt no later than the second{-}to{-}last turn.} \\
\texttt{Do not wait until the very end.} \\
\\ 
\texttt{You can use these actions:} \\
\texttt{{-} JSON format described at the end of the prompt} \\
\\ 
\\ 
\texttt{Do not add import statements in your guess.} \\
\texttt{Anything you might need is already available, so just write the function body in hypothesis and return the answer.} \\
\\ 
\texttt{Preloaded examples:} \\
\texttt{1. ({[}5, {-}5, 5{]}, 3)} \\
\texttt{2. ({[}0{]}, 0)} \\
\texttt{3. ({[}1, 1{]}, 2)} \\
\texttt{4. ({[}3, 2, 1{]}, 2)} \\
\texttt{5. ({[}1, 2, 3{]}, 2)} \\
\texttt{6. ({[}2, 1{]}, 1)} \\
\texttt{7. ({[}{]}, 0)} \\
\texttt{8. ({[}{-}1, 0, 1{]}, 2)} \\
\texttt{9. ({[}2{]}, 0)} \\
\texttt{10. ({[}2, 18, {-}1{]}, 1)} \\
\\ 
\texttt{Scoring:} \\
\texttt{{-} correctness is required for any non{-}zero score} \\
\texttt{{-} in interactive modes, solving with fewer information requests is better} \\
\\ 
\texttt{Input formatting reminders:} \\
\texttt{{-} strings must be quoted} \\
\texttt{{-} booleans must be True/False} \\
\texttt{{-} use commas between multiple values} \\
\texttt{{-} do not write any extra action tags in the same message} \\
\\ 
\texttt{Format example:} \\
\texttt{You write:} \\
\texttt{\{'current\_mode': <probing/confirming/don't know; choose whichever best describes your current state in the task>, 'input\_rationale': <brief description of why you chose this hypothesis>, 'hypothesis': <either the exact string "unknown" if you have no guess yet OR your current best guess as runnable Python matching the given signature, with a one{-}line docstring stating what it does. Never write placeholder code, if you cannot commit to a guess, use "unknown". Match the given signature.} \\
\texttt{e.g.} \\
\texttt{def f(x):} \\
\texttt{	"""Returns the absolute value of x."""} \\
\texttt{	return x if x >= 0 else {-}x} \\
\\ 
\texttt{Use "unknown" if you have no guess.>, 'confidence\_score': <score between 1{-}100. 100 means highest confidence in your hypothesis>\}} \\
\\ 
\texttt{Output only one valid json per turn.} \\
            }
        }
    }
     \\ \\

    \theutterance \stepcounter{utterance}  
    & \multicolumn{4}{p{0.6\linewidth}}{
        \cellcolor[rgb]{0.561,0.737,0.561}{
            \makecell[{{p{\linewidth}}}]{
                \texttt{\tiny{[P1$\rangle$GM]}}
                \texttt{```json} \\
\texttt{\{} \\
\texttt{  "current\_mode": "confirming",} \\
\texttt{  "input\_rationale": "Analyzing the examples, the function seems to count the number of elements that are greater than all previous elements in the list. For example, in {[}5, {-}5, 5{]}, 5 is greater than nothing (count 1), {-}5 is not greater than 5, and 5 is greater than {-}5 (count 2), total 2, but the output is 3, which doesn't match. Alternatively, it might count the number of elements that are greater than or equal to the current maximum at their position. However, the first example ({[}5, {-}5, 5{]}) outputs 3, which suggests it counts all elements that are not less than the first element (5, {-}5 is less, 5 is equal, so 2), but the output is 3, which contradicts. Another observation: in {[}3, 2, 1{]}, output is 2, which could be the count of elements that are not less than the last element (3 and 2 are >=1, so 2). Similarly, {[}1, 2, 3{]} outputs 2, which could be the count of elements >= the first element (1 and 2, but 3 is also >=1, so 3, which contradicts). The pattern is unclear. Another guess: it counts the number of elements that} \\
            }
        }
    }
    & & \\ \\

    \theutterance \stepcounter{utterance}  
    & & & \multicolumn{2}{p{0.3\linewidth}}{
        \cellcolor[rgb]{0.9,0.9,0.9}{
            \makecell[{{p{\linewidth}}}]{
                \texttt{\tiny{[GM$|$GM]}}
                \texttt{Game Verdict: LOSS (turn limit reached)} \\
            }
        }
    }
    & & \\ \\

\end{supertabular}
}

\end{document}
